\documentclass[11pt]{article}

\usepackage[T1]{fontenc}
\usepackage{amsmath,amssymb,amsthm}
\usepackage[margin=1in]{geometry}
\usepackage[hidelinks]{hyperref}
\usepackage{booktabs}

\newtheorem{theorem}{Theorem}
\newtheorem{lemma}[theorem]{Lemma}
\newtheorem{remark}[theorem]{Remark}

% Math alphabets are written explicitly below for deterministic citation parsing.
% No user-defined shorthand commands are needed.

\title{A Degree-Four Obstruction to Nonnegative Closure for Independent-Row $K$-Species}
\author{Xiaojian Zeng\\
PBC School of Finance, Tsinghua University, 100190\\
\texttt{xiaojianzeng@foxmail.com}}
\date{}

\begin{document}
\maketitle

\begin{abstract}
For a partition $\alpha=(a_1,\ldots,a_k)$ of $n$, let $G_\alpha$
be the direct product of the cyclic groups generated by the disjoint row
cycles of lengths $a_1,\ldots,a_k$, and put
$\mathbf{K}_\alpha=X^n/G_\alpha$.  We prove that every Hadamard product
$\mathbf{K}_\alpha\times\mathbf{K}_\beta$ has a nonnegative integral decomposition into
independent-row $K$-species in degrees $1$, $2$, and $3$, but that this
property fails in degree $4$.  More precisely,
$\mathbf{K}_{(2,2)}\times\mathbf{K}_{(4)}$ has two molecular orbits, one of whose
stabilizers is generated by a double transposition, and its unique
degree-four cycle-index expansion is
\[
  2Z_{\mathbf{K}_{(2,2)}}-2Z_{\mathbf{K}_{(2,1,1)}}+2Z_{\mathbf{K}_{(1,1,1,1)}}.
\]
Thus degree four is the smallest degree of nonnegative-closure failure for
this convention.  The result does not address closed formulas in a larger
molecular-species basis or the corresponding problem under a different
repeated-part convention.
\end{abstract}

\section{Introduction and convention}

Combinatorial species and their cycle indices provide a functorial form of
labelled enumeration \cite{Joyal1981,BergeronLabelleLeroux1998}.  The
internal, or Kronecker, product of symmetric functions is reviewed in
\cite[Chapter~I]{Macdonald1995}; in the species setting it is induced by the
pointwise Cartesian product.  Positive expansions for particular bases are
therefore naturally tied to orbit stabilizers.  The classical complete
homogeneous basis has such a combinatorial product rule
\cite{GarsiaRemmel1985}.

Two recent papers motivate the precise convention studied here.  In
\cite[Definition~3.3]{BaolahyBenjamin2026}, equal rows of a fixed length are
rotated diagonally by one cyclic factor, and positive closure is proved for
the resulting species.  By contrast, \cite[Remark~3.2 and Open
Problem~5.1]{BaolahyRandrianirina2026} explicitly introduces a companion
group with one independent cyclic factor per row.  The conventions may
differ when a nontrivial part is repeated: for example, the first gives
$\langle(12)(34)\rangle$ at shape $(2,2)$, whereas the second gives
$\langle(12),(34)\rangle$.  The present note concerns only the second,
independent-row convention.  Consequently it neither contradicts the
positive theorem proved for the first convention nor settles every part of
the broader closed-form problem posed in \cite{BaolahyRandrianirina2026}.

Coefficients in a specified symmetric-function basis can encode sharply
delimited combinatorial information in other settings as well.  For example,
leading and adjacent star-basis layers of the chromatic symmetric function
recover weighted-core data for a class of trees
\cite{Zeng2026MultiplicityArxiv}.  That invariant is different from the
cycle indices below; the example is included only to situate the
basis-sensitive method.

Let $n\geq1$ and let $\alpha=(a_1,\ldots,a_k)$ be a partition of $n$.  On consecutive,
disjoint blocks of $[n]$, let $c_i$ be the $a_i$-cycle (with a $1$-cycle
understood as the identity), and define
\[
 G_\alpha=\langle c_1\rangle\times\cdots\times\langle c_k\rangle
 \leq \mathfrak{S}_n,
 \qquad
 \mathbf{K}_\alpha=X^n/G_\alpha.
\]
For $H\leq\mathfrak{S}_n$, the molecular species $X^n/H$ has, on an $n$-element
label set, the transitive $\mathfrak{S}_n$-set $\mathfrak{S}_n/H$.  Its cycle index is
\begin{equation}
 Z_{X^n/H}=\frac{1}{|H|}\sum_{h\in H}p_{\operatorname{type}(h)}.
 \label{eq:cycle-index}
\end{equation}
The pointwise Cartesian product of species is denoted by $\times$, and its
cycle index obeys
\begin{equation}
 Z_{F\times G}=Z_F\star Z_G,
 \qquad
 p_\lambda\star p_\mu=\delta_{\lambda\mu}z_\lambda p_\lambda,
 \label{eq:internal-product}
\end{equation}
where $z_\lambda=\prod_i i^{m_i}m_i!$ for
$\lambda=(1^{m_1}2^{m_2}\cdots)$.

Our main result resolves nonnegative closure, not the existence of an
expression over arbitrary molecular species.

\begin{theorem}\label{thm:main}
Under the independent-row convention above, for every $n\in\{1,2,3\}$ and
every pair of partitions $\alpha,\beta$ of $n$, the species
$\mathbf{K}_\alpha\times\mathbf{K}_\beta$ is a coproduct of the $\mathbf{K}_\lambda$ with
nonnegative integral multiplicities.  In degree four,
$\mathbf{K}_{(2,2)}\times\mathbf{K}_{(4)}$ admits no such decomposition.  Its unique
expansion at the level of cycle indices is
\begin{equation}
 Z_{\mathbf{K}_{(2,2)}}\star Z_{\mathbf{K}_{(4)}}
 =2Z_{\mathbf{K}_{(2,2)}}-2Z_{\mathbf{K}_{(2,1,1)}}
   +2Z_{\mathbf{K}_{(1,1,1,1)}}.
 \label{eq:signed-expansion}
\end{equation}
Hence degree four is the smallest degree in which nonnegative closure fails.
\end{theorem}

\section{Double cosets and the degree-four obstruction}

We first record the orbit decomposition in the exact form needed below.

\begin{lemma}[Coset-product lemma]\label{lem:coset}
For $H,L\leq\mathfrak{S}_n$ there is a species isomorphism
\begin{equation}
 (X^n/H)\times(X^n/L)
 \cong
 \coprod_{r\in H\backslash\mathfrak{S}_n/L}
 X^n/(H\cap rLr^{-1}).
 \label{eq:coset-decomposition}
\end{equation}
Moreover,
\begin{equation}
 |HrL|=\frac{|H||L|}{|H\cap rLr^{-1}|}.
 \label{eq:double-coset-size}
\end{equation}
\end{lemma}

\begin{proof}
On the standard label set the diagonal $\mathfrak{S}_n$-action on
$\mathfrak{S}_n/H\times\mathfrak{S}_n/L$ has representatives $(H,rL)$.  Two representatives
belong to the same orbit exactly when their parameters lie in the same
double coset $HrL$, and the stabilizer of $(H,rL)$ is
$H\cap rLr^{-1}$.  Decomposing into transitive orbits gives (4).
Equation (5) follows by counting the fibers of the multiplication map
$H\times L\to HrL$, $(h,\ell)\mapsto hr\ell$.
\end{proof}

For the target pair put
\[
 H=G_{(2,2)}=\langle(12),(34)\rangle,
 \qquad L=G_{(4)}=\langle(1234)\rangle,
 \qquad \pi=(23).
\]
The groups are
\[
 H=\{e,(12),(34),(12)(34)\},
 \quad
 L=\{e,(1234),(13)(24),(1432)\}.
\]
Thus $H\cap L=\{e\}$ and $|HL|=16$.  Conjugating the generator of $L$ by
$\pi$ gives $(1324)$, whose square is $(12)(34)$.  Its two elements of
order four do not lie in $H$, and hence
\[
 H\cap\pi L\pi^{-1}=\{e,(12)(34)\},
 \qquad |H\pi L|=8.
\]
The two double cosets are distinct and their sizes sum to $24$; they
therefore exhaust $\mathfrak{S}_4$.  Lemma~\ref{lem:coset} yields the more precise
molecular decomposition
\begin{equation}
 \mathbf{K}_{(2,2)}\times\mathbf{K}_{(4)}
 \cong \mathbf{K}_{(1,1,1,1)}\;\coprod\;
 X^4/\langle(12)(34)\rangle.
 \label{eq:molecular-decomposition}
\end{equation}
The second stabilizer is not conjugate to any independent-row $G_\lambda$.
Indeed, the group orders are
\[
\begin{array}{c|ccccc}
\lambda&(4)&(3,1)&(2,2)&(2,1,1)&(1^4)\\ \hline
|G_\lambda|&4&3&4&2&1.
\end{array}
\]
The only order-two candidate is generated by a transposition, which is not
conjugate to a double transposition.

To turn this stabilizer obstruction into a proof that no nonnegative
$K$-decomposition exists, we establish uniqueness of the cycle-index
coordinates.  In the power-sum order
$(p_4,p_{31},p_{22},p_{211},p_{1^4})$, direct averaging in (1) gives
\begin{align}
 Z_{\mathbf{K}_4}&=\tfrac14(p_{1^4}+p_{22}+2p_4),&
 Z_{\mathbf{K}_{31}}&=\tfrac13(p_{1^4}+2p_{31}),\notag\\
 Z_{\mathbf{K}_{22}}&=\tfrac14(p_{1^4}+2p_{211}+p_{22}),&
 Z_{\mathbf{K}_{211}}&=\tfrac12(p_{1^4}+p_{211}),\notag\\
 Z_{\mathbf{K}_{1^4}}&=p_{1^4}. \label{eq:degree-four-indices}
\end{align}
The matrix having these vectors as columns, in partition order
$(4),(31),(22),(211),(1^4)$, is
\[
 \begin{pmatrix}
 1/2&0&0&0&0\\
 0&2/3&0&0&0\\
 1/4&0&1/4&0&0\\
 0&0&1/2&1/2&0\\
 1/4&1/3&1/4&1/2&1
 \end{pmatrix}.
\]
Its determinant is $1/24$, so the five degree-four cycle indices are
linearly independent.

The only cycle types common to $Z_{\mathbf{K}_{22}}$ and $Z_{\mathbf{K}_4}$ are $1^4$ and
$22$.  Since $z_{1^4}=24$ and $z_{22}=8$, (2) gives
\begin{equation}
 Z_{\mathbf{K}_{22}}\star Z_{\mathbf{K}_4}
 =\tfrac32p_{1^4}+\tfrac12p_{22}. \label{eq:target-index}
\end{equation}
Comparing successively the coefficients of
$p_4,p_{31},p_{22},p_{211},p_{1^4}$ in (7)--(8) forces the coordinate vector
\[
 (c_4,c_{31},c_{22},c_{211},c_{1^4})=(0,0,2,-2,2),
\]
which is (3).  Linear independence makes this expansion unique.  A species
coproduct with nonnegative multiplicities would induce a nonnegative
cycle-index expansion, contrary to the coefficient $-2$.  Notice that this
argument uses only the forward functoriality of the cycle index; it does not
assume that equality of arbitrary cycle indices implies species isomorphism.

\section{Exhaustion of the lower degrees}

It remains to prove the minimal-degree assertion.  Lemma~\ref{lem:coset}
gives all products directly.  In degree one,
\[
 \mathbf{K}_1\times\mathbf{K}_1\cong\mathbf{K}_1.
\]
In degree two the complete table of unordered products is
\[
\begin{array}{c|cc}
\times&\mathbf{K}_2&\mathbf{K}_{11}\\ \hline
\mathbf{K}_2&\mathbf{K}_2&\mathbf{K}_{11}\\
\mathbf{K}_{11}&&2\mathbf{K}_{11}
\end{array}
\]
because $G_2=\mathfrak{S}_2$ and $G_{11}=1$.

For degree three write $A=G_3=A_3$, let $B=G_{21}$ be a transposition
subgroup, and let $T=G_{111}=1$.  The complete table is
\begin{equation}
\begin{array}{c|ccc}
\times&\mathbf{K}_3&\mathbf{K}_{21}&\mathbf{K}_{111}\\ \hline
\mathbf{K}_3&2\mathbf{K}_3&\mathbf{K}_{111}&2\mathbf{K}_{111}\\
\mathbf{K}_{21}&&\mathbf{K}_{21}+\mathbf{K}_{111}&3\mathbf{K}_{111}\\
\mathbf{K}_{111}&&&6\mathbf{K}_{111}
\end{array}. \label{eq:degree-three-table}
\end{equation}
Here $A$ is normal of index two, so the $(A,A)$ product has two orbits, both
with stabilizer $A$.  Every conjugate of $B$ meets $A$ trivially, and
$|A||B|=6$, so the $(A,B)$ product has one free orbit.  For $(B,B)$, the
identity double coset has stabilizer $B$ and size two; the remaining double
coset has trivial stabilizer and size four.  Finally, pairing any subgroup
with $T$ gives its index many free orbits.  These arguments prove every
entry of (9) as a species isomorphism.  The displayed tables exhaust all
partition pairs in degrees below four and complete the proof of
Theorem~\ref{thm:main}.

\section{Exact computational corroboration}

The proof above is self-contained.  It is also corroborated by two exact
implementations supplied with the reproducibility archive.  The first
constructs every independent-row group and every double coset for all $39$
ordered partition pairs in degrees $1$ through $4$, records the stabilizer
intersections, computes cycle indices over $\mathbb Q$, and solves every
basis system.  A second implementation imports no producer code: it builds
row groups as Cartesian products of rotations, finds double cosets as
$H\times L$-orbits by breadth-first search, and solves the basis systems by
Leibniz determinants and Cramer's rule.  It independently checks $167$
double cosets and $663$ double-coset element memberships, recovers the
degree-wise determinants
\[
 1,\quad \tfrac12,\quad \tfrac13,\quad \tfrac1{24},
\]
and returns $(0,0,2,-2,2)$ for the target pair.  Both implementations use
integer arithmetic and exact rational arithmetic only; these computations
are corroboration rather than premises replacing the proof.

\section{Limitations and assistance disclosure}

Equation (6) already supplies a positive molecular decomposition if one
allows the additional molecule $X^4/\langle(12)(34)\rangle$.  The theorem
therefore says specifically that the independent-row family
$\{\mathbf{K}_\lambda:\lambda\text{ is a partition of }n\}$ is not nonnegatively closed.  It does not
exclude a closed formula in a larger molecular basis, does not address the
diagonal repeated-part convention of \cite{BaolahyBenjamin2026}, and does
not resolve the remaining directions in Open Problem~5.1 of
\cite{BaolahyRandrianirina2026}.

OpenAI Codex assisted with organizing the audited argument, bibliography
triage, and preparation of the \LaTeX{} manuscript.  The exact computations
were performed by the two deterministic programs described above.  No
floating-point calculation is used in the mathematical conclusion.  The
author reviewed the proof, citations, and computational claims and assumes
responsibility for the final text.

\bibliographystyle{plain}
\bibliography{references}

\end{document}
